\chapter{Liver Cirrhosis}
\index{Liver Cirrhosis}

\section*{Definition and Overview}
Liver cirrhosis is the late, irreversible stage of progressive hepatic fibrosis characterised by the distortion of normal liver architecture. Chronic inflammation and recurrent hepatocyte injury stimulate the deposition of extracellular matrix (collagen), resulting in diffuse fibrotic scar tissue that surrounds regenerative nodules of hepatocytes. This structural reorganization obstructs portal blood flow and compromises hepatic cellular function, culminating in portal hypertension and progressive liver failure. Cirrhosis is clinically divided into compensated and decompensated phases, the latter being defined by the development of life-threatening complications such as ascites, hepatic encephalopathy, and variceal haemorrhage.

\section*{Epidemiology}
Globally, liver cirrhosis represents a major public health burden, causing approximately 1.3 million deaths annually and ranking as a leading cause of mortality. In many countries and other developed countries, the epidemiology of cirrhosis is evolving. While alcohol-related liver disease and chronic Hepatitis C have historically been the primary drivers, metabolic dysfunction-associated steatotic liver disease (MASLD, formerly known as non-alcoholic fatty liver disease) has emerged as a major cause due to rising obesity and diabetes rates. Cirrhosis is more prevalent in males and is most commonly diagnosed in individuals aged between 50 and 70 years.

\section*{Aetiology and Pathophysiology}
Cirrhosis develops from prolonged hepatic insult from various sources, triggering a fibrogenic cascade mediated by hepatic stellate cells. Aetiological factors and their pathophysiology include:

\begin{itemize}
    \item \textbf{Chronic alcohol misuse:} Direct metabolic toxicity to hepatocytes, leading to intracellular fat accumulation (steatosis), lipid peroxidation, inflammation, and eventual collagen deposition.
    \item \textbf{Metabolic dysfunction-associated steatohepatitis (MASH):} Hepatic fat accumulation associated with insulin resistance and obesity, which triggers chronic lipotoxicity and inflammatory activation.
    \item \textbf{Chronic viral hepatitis:} Persistent infection with Hepatitis B virus (HBV) or Hepatitis C virus (HCV) inducing immune-mediated cytolytic damage and chronic fibrogenesis.
    \item \textbf{Autoimmune liver diseases:} Targeted destruction of hepatic tissue by the immune system, including autoimmune hepatitis, primary biliary cholangitis (PBC), and primary sclerosing cholangitis (PSC).
    \item \textbf{Genetic and metabolic disorders:} Hereditary haemochromatosis (excessive iron deposition), Wilson's disease (impaired copper excretion), and alpha-1 antitrypsin deficiency (protein misfolding causing hepatocyte death).
\end{itemize}

\section*{Clinical Features}
The clinical presentation of cirrhosis ranges from complete absence of symptoms in the compensated phase to severe multi-system dysfunction in the decompensated phase. Features include:

\begin{itemize}
    \item \textbf{Compensated symptoms:} Non-specific features such as chronic fatigue, lethargy, generalised malaise, mild right upper quadrant abdominal discomfort, and unexplained weight loss.
    \item \textbf{Decompensated signs:} Jaundice (yellowing of the skin and sclera due to hyperbilirubinaemia), ascites (tense abdominal swelling from fluid accumulation), and peripheral pitting oedema.
    \item \textbf{Hepatic encephalopathy:} Cognitive fluctuations, sleep-wake cycle reversal, impaired concentration, confusion, slurred speech, and asterixis (a characteristic flapping tremor of the hands).
    \item \textbf{Gastrointestinal haemorrhage:} Haematemesis (vomiting blood) or melaena (passing black, tarry stools) resulting from ruptured gastro-oesophageal varices.
    \item \textbf{Dermatological signs:} Spider naevi (telangiectasias on the upper chest/arms), palmar erythema (redness of the palms), caput medusae (dilated periumbilical veins), pruritus, and petechiae or easy bruising.
    \item \textbf{Endocrine and nutritional features:} Muscle wasting (sarcopenia), testicular atrophy, gynaecomastia, and loss of secondary sexual hair in males due to impaired hepatic metabolism of oestrogens.
\end{itemize}

\section*{Differential Diagnosis}
When presenting with signs of portal hypertension or hepatic dysfunction, several other conditions must be considered:

\begin{itemize}
    \item \textbf{Acute liver failure:} Severe hepatic dysfunction and encephalopathy occurring within weeks of an insult (e.g.\ paracetamol overdose) in a patient without underlying chronic liver disease.
    \item \textbf{Non-cirrhotic portal hypertension:} Portal hypertension caused by conditions such as portal vein thrombosis, splenic vein thrombosis, or schistosomiasis, without the histopathological changes of cirrhosis.
    \item \textbf{Congestive cardiac failure:} Severe right-sided heart failure causing passive hepatic congestion (cardiac cirrhosis), presenting with hepatomegaly, ascites, and elevated jugular venous pressure.
    \item \textbf{Budd-Chiari syndrome:} Obstruction of hepatic venous outflow (e.g.\ due to hepatic vein thrombosis), presenting with acute abdominal pain, hepatomegaly, and rapid-onset ascites.
    \item \textbf{Infiltrative liver disease:} Amyloidosis, sarcoidosis, or diffuse metastatic malignancy that presents with hepatomegaly and abnormal liver biochemistry mimicking cirrhosis.
\end{itemize}

\section*{When to Seek Medical Care}
Patients with a history of liver disease require regular medical monitoring. However, certain symptoms indicate acute clinical deterioration or life-threatening complications.

\textbf{Contact your local emergency services for:}
\begin{itemize}
    \item Vomiting blood (haematemesis), coffee-ground emesis, or passing black, tarry stools (melaena), which indicate active, life-threatening variceal haemorrhage.
    \item Acute onset of severe confusion, extreme drowsiness, slurred speech, or coma (suggesting advanced hepatic encephalopathy).
    \item Sudden, severe abdominal pain or a high fever with chills, which can indicate spontaneous bacterial peritonitis.
    \item Progressive shortness of breath or chest pain, which may occur with tense ascites or hepatic hydrothorax.
\end{itemize}

\section*{Diagnosis and Investigations}
A combination of laboratory tests, imaging, and endoscopic evaluations is used to diagnose cirrhosis, assess its severity, and screen for complications. Investigations include:

\begin{itemize}
    \item \textbf{Liver function tests (LFTs):} Typically demonstrating elevated serum bilirubin, alkaline phosphatase (ALP), gamma-glutamyl transferase (GGT), and aspartate/alanine aminotransferases (AST/ALT), alongside decreased serum albumin.
    \item \textbf{Coagulation studies:} Prolonged prothrombin time (PT) and elevated International Normalised Ratio (INR), reflecting impaired hepatic synthesis of vitamin K-dependent clotting factors.
    \item \textbf{Full blood count (FBC):} Thrombocytopenia (low platelet count) is a highly sensitive non-invasive indicator of portal hypertension and subsequent splenic sequestration (hypersplenism).
    \item \textbf{Transient elastography (FibroScan) and ultrasound:} Abdominal ultrasound assesses liver morphology, portal vein flow, and ascites; FibroScan measures liver stiffness to non-invasively quantify the degree of fibrosis.
    \item \textbf{Upper endoscopy (gastroscopy):} Essential to screen for the presence of gastro-oesophageal varices and determine the risk of bleeding.
    \item \textbf{Liver biopsy:} The historical gold standard for diagnosis, now reserved for cases where the diagnostic etiology remains uncertain despite non-invasive testing.
\end{itemize}

\section*{Management}
Management is directed at halting disease progression, managing portal hypertension, treating decompensated complications, and assessing suitability for transplantation. Interventions include:

\begin{itemize}
    \item \textbf{Etiology-specific therapies:} Strict abstinence from alcohol, antiviral suppression for Hepatitis B/C, weight loss and metabolic control for MASH, or corticosteroid therapy for autoimmune hepatitis.
    \item \textbf{Portal hypertension management:} Use of non-selective beta-blockers (e.g.\ carvedilol or propranolol) or endoscopic variceal ligation (EVL) to reduce the risk of variceal bleeding.
    \item \textbf{Ascites control:} A combination of dietary sodium restriction (less than 2 g/day) and oral diuretics (spironolactone and furosemide), with therapeutic large-volume paracentesis (plus albumin infusion) for refractory ascites.
    \item \textbf{Encephalopathy therapy:} Oral lactulose to promote bowel movements and trap ammonia in the gut, often combined with rifaximin (a non-absorbable antibiotic) to reduce ammonia-producing intestinal flora.
    \item \textbf{Nutritional optimization:} Providing a high-calorie, high-protein diet (1.2--1.5 g/kg/day of protein) and small, frequent meals (including a late-night snack) to prevent muscle wasting.
    \item \textbf{Liver transplantation:} The only curative treatment, indicated for patients with decompensated cirrhosis (assessed via the MELD score) or those with early-stage hepatocellular carcinoma.
\end{itemize}

\section*{Complications}
The progression of cirrhosis leads to severe systemic complications, including:

\begin{itemize}
    \item \textbf{Variceal haemorrhage:} Life-threatening upper gastrointestinal bleeding from ruptured oesophageal or gastric varices.
    \item \textbf{Spontaneous bacterial peritonitis (SBP):} A bacterial infection of the ascitic fluid, requiring urgent intravenous antibiotics.
    \item \textbf{Hepatorenal syndrome (HRS):} Functional renal failure driven by severe renal vasoconstriction secondary to splanchnic vasodilation.
    \item \textbf{Hepatocellular carcinoma (HCC):} Primary liver cancer; patients with cirrhosis have an annual risk of 1--8\% and require lifelong 6-monthly ultrasound surveillance.
    \item \textbf{Hepato-pulmonary syndrome:} Hypoxaemia caused by intrapulmonary vascular dilations in the setting of portal hypertension.
\end{itemize}

\section*{Prognosis and Outcomes}
The prognosis of liver cirrhosis depends on whether the patient is compensated or decompensated. Patients with compensated cirrhosis have a median survival exceeding 12 years. Once a decompensating event occurs, the median survival drops precipitously to approximately 1.5 to 2 years. Prognosis is clinically calculated using the Child-Pugh classification (based on bilirubin, albumin, INR, ascites, and encephalopathy) and the Model for End-Stage Liver Disease (MELD) score, which predicts 3-month mortality and guides liver transplant allocation.

\section*{Prevention and Self-Care}
Patients can take active measures to slow the progression of cirrhosis and avoid acute decompensation:

\begin{itemize}
    \item \textbf{Strict alcohol avoidance:} Essential for all patients, regardless of the primary cause of their cirrhosis, to prevent accelerated liver injury.
    \item \textbf{Strict low-sodium diet:} Limiting sodium intake to prevent fluid retention, ascites, and peripheral swelling.
    \item \textbf{Avoid hepatotoxic substances:} Avoiding non-steroidal anti-inflammatory drugs (NSAIDs, which increase the risk of renal failure and variceal bleeding) and keeping paracetamol intake to a safe level (maximum 2 g/day).
    \item \textbf{Receive routine vaccinations:} Ensuring vaccination against Hepatitis A, Hepatitis B, influenza, and pneumococcal disease to prevent severe infections.
    \item \textbf{Adhere to surveillance programmes:} Attending all scheduled 6-monthly ultrasound scans for hepatocellular carcinoma screening and routine diagnostic blood tests.
\end{itemize}

\section*{Sources and Further Reading}
The following organisations provide reputable details and clinical guidelines on liver cirrhosis:

\begin{itemize}
    \item Gastroenterological Society of Australia. \textit{Clinical resources and patient guides on liver cirrhosis and portal hypertension.} https://www.gesa.org.au/
    \item Healthdirect Australia. \textit{Consumer guide on liver cirrhosis, explaining symptoms, causes, and treatment options.} https://www.healthdirect.gov.au/
    \item Mayo Clinic. \textit{Comprehensive medical information on liver cirrhosis, diagnostic criteria, and management.} https://www.mayoclinic.org/
    \item NHS. \textit{Patient resource on cirrhosis, covering causes, complications, and lifestyle advice.} https://www.nhs.uk/
    \item World Health Organization. \textit{Global reports on the burden of hepatitis and cirrhosis, and prevention strategies.} https://www.who.int/
\end{itemize}
